% Luc Destombe --- licence CC BY-NC-SA 4.0
% https://creativecommons.org/licenses/by-nc-sa/4.0/deed.fr
% Fichier autonome : compiler avec pdflatex, deux fois (signets, nombre de pages).
\documentclass[10pt,a4paper,twocolumn]{article}
\makeatletter
\newif\iflycee@unecolonne
\newif\iflycee@palatino
\lycee@palatinotrue

\RequirePackage[utf8]{inputenc}
\RequirePackage[T1]{fontenc}
\RequirePackage[french]{babel}
\RequirePackage{etoolbox}

\newcommand{\ShorthandsOff}{\@ifundefined{shorthandoff}{}{\shorthandoff{;:!?}}}
\newcommand{\ShorthandsOn}{\@ifundefined{shorthandon}{}{\shorthandon{;:!?}}}
\ShorthandsOff
\AtBeginDocument{\ShorthandsOn}
\AtBeginEnvironment{tikzpicture}{\ShorthandsOff}

\RequirePackage{geometry}
\RequirePackage{fancyhdr}
\RequirePackage{lastpage}
\RequirePackage{multicol}
\RequirePackage{array}
\RequirePackage{enumitem}
\RequirePackage{xcolor}
\RequirePackage{needspace}

\RequirePackage{amsmath,amssymb}
\RequirePackage{mathpazo}
\DeclareMathSizes{10.95}{10.95}{8}{5.5}
\begingroup\catcode`\_=\active \catcode`\^=\active
\gdef_#1{\sb{\scriptscriptstyle #1}}
\gdef^#1{\sp{\scriptscriptstyle\lycee@moinsepais #1}}
\endgroup
\newcommand\lycee@moins{\mathbin{%
  \setbox\z@\hbox{$\m@th\scriptscriptstyle\lycee@moinsorig$}%
  \dimen@=.55\wd\z@ \advance\dimen@-.5pt
  \hbox to.75\wd\z@{\hss$\m@th\scriptscriptstyle\vcenter{\hbox{%
    \tikz\draw[line width=.5pt,line cap=round](0,0)--(\the\dimen@,0);}}$\hss}}}
\begingroup\lccode`\~=`\- \lowercase{\endgroup\let~}\lycee@moins
\newcommand\lycee@moinsepais{\mathcode`\-="8000 }
\AtBeginDocument{\mathchardef\lycee@moinsorig=\mathcode`\-\relax}
\AtBeginDocument{\catcode`\_=12 \mathcode`\_="8000
  \catcode`\^=12 \mathcode`\^="8000 }
\newcommand\lycee@exposants{%
  \fontdimen13\textfont2=.48\fontdimen6\textfont2
  \fontdimen14\textfont2=.48\fontdimen6\textfont2
  \fontdimen15\textfont2=.48\fontdimen6\textfont2
  \fontdimen13\scriptfont2=.48\fontdimen6\scriptfont2
  \fontdimen14\scriptfont2=.48\fontdimen6\scriptfont2
  \fontdimen15\scriptfont2=.48\fontdimen6\scriptfont2}
\every@math@size\expandafter{\the\every@math@size\lycee@exposants}
\newlength{\retraitcalculs}
\setlength{\retraitcalculs}{2em}

\RequirePackage{tikz}
\usetikzlibrary{arrows.meta,calc,babel,shapes.misc}
\RequirePackage[most]{tcolorbox}
\tcbuselibrary{skins,breakable}

\definecolor{coulDef}{HTML}{1F4E79}
\definecolor{coulProp}{HTML}{7030A0}
\definecolor{coulRem}{HTML}{C55A11}
\definecolor{coulEx}{HTML}{2E7D32}
\definecolor{coulExo}{HTML}{B03A2E}
\definecolor{coulPart}{HTML}{1F4E79}
\definecolor{coulMeth}{HTML}{1B6E4F}
\definecolor{coulCourbe}{HTML}{0B5ED7}
\definecolor{coulCourbeB}{HTML}{C00000}
\definecolor{coulCourbeC}{HTML}{2E7D32}
\definecolor{coulGrille}{HTML}{8C8C8C}
\definecolor{coulRep}{HTML}{1B6E4F}
\definecolor{coulBar}{HTML}{2E7D32}

\newcommand{\R}{\mathbb{R}}
\newcommand{\Rs}{\mathbb{R}^{*}}
\renewcommand{\le}{\leqslant}
\renewcommand{\ge}{\geqslant}
\newcommand{\vect}[1]{\overrightarrow{#1}}
\newcommand{\norme}[1]{\left\lVert #1 \right\rVert}
\newcommand{\intervalle}[2]{\left[#1\,;#2\right]}
\RequirePackage{cancel}
\renewcommand{\CancelColor}{\color{coulExo}}
\newcommand{\fois}[1]{\times{\color{coulExo}#1}}
\newcommand{\barre}[1]{\cancel{#1}}

\tikzset{grille/.style={coulGrille,line width=0.4pt}}
\newcommand{\quadrillage}[4]{\draw[grille,step=1] (#1,#2) grid (#3,#4);}
\tikzset{
  croix/.style={cross out,draw,line width=0.6pt,inner sep=0pt,minimum size=4pt},
  croixdroite/.style={croix,rotate=45,line width=0.8pt},
}
\newcommand{\pt}[3]{\path (#1) node[croix]{} node[#2,font=\small,inner sep=2.5pt]{$#3$};}

\newcommand{\lycee@repere}[4]{%
  \draw[grille,step=1] (#1,#2) grid (#3,#4);
  \draw[-{Stealth[length=2mm]},thick] (#1,0)--({#3+0.4},0) node[below left=-1pt]{$x$};
  \draw[-{Stealth[length=2mm]},thick] (0,#2)--(0,{#4+0.4}) node[below left=-1pt]{$y$};
  \node[below left,font=\scriptsize,inner sep=1.5pt] at (0,0) {$0$};}
\newcommand{\lycee@gradx}[1]{\foreach \k/\l in {#1}{%
  \draw (\k,0.07)--(\k,-0.07) node[below,font=\scriptsize,inner sep=1.5pt]{$\l$};}}
\newcommand{\lycee@grady}[1]{\foreach \k/\l in {#1}{%
  \draw (0.07,\k)--(-0.07,\k) node[left,font=\scriptsize,inner sep=1.5pt]{$\l$};}}
\newcommand{\lycee@intff}[2]{\left[#1\,;#2\right]}
\newcommand{\lycee@intfo}[2]{\left[#1\,;#2\right[}
\newcommand{\lycee@intof}[2]{\left]#1\,;#2\right]}
\newcommand{\lycee@intoo}[2]{\left]#1\,;#2\right[}
\newcommand{\lycee@ens}[1]{\left\{#1\right\}}
\AtBeginDocument{%
  \providecommand{\repere}{\lycee@repere}%
  \providecommand{\gradx}{\lycee@gradx}%
  \providecommand{\grady}{\lycee@grady}%
  \providecommand{\Cf}{\mathcal{C}\sb{\scriptscriptstyle f}}%
  \providecommand{\Cg}{\mathcal{C}\sb{\scriptscriptstyle g}}%
  \providecommand{\intff}{\lycee@intff}%
  \providecommand{\intfo}{\lycee@intfo}%
  \providecommand{\intof}{\lycee@intof}%
  \providecommand{\intoo}{\lycee@intoo}%
  \providecommand{\ens}{\lycee@ens}}

\newlist{questions}{enumerate}{3}
\setlist[questions,1]{label=\arabic*\textdegree),leftmargin=*,itemsep=2pt,topsep=3pt}
\setlist[questions,2]{label=\alph*),leftmargin=*,itemsep=1pt,topsep=2pt}
\setlist[questions,3]{label=\roman*),leftmargin=*,itemsep=1pt,topsep=1pt}
\setlist[itemize]{label=\textbullet,leftmargin=*,itemsep=2pt,topsep=3pt}

\newcommand{\besoinplace}[1]{\par\penalty\@M\begingroup
  \dimen@=#1\relax\dimen@ii=\pagegoal\advance\dimen@ii-\pagetotal
  \ifdim\dimen@>\dimen@ii\ifdim\dimen@ii>\z@\vfil\fi\break\fi\endgroup}

\newcommand{\lycee@pied}{\thepage/\pageref{LastPage}}
\newcommand{\entete}[2]{%
  \pagestyle{fancy}\fancyhf{}%
  \fancyhead[L]{\footnotesize\color{coulPart}\textbf{#1}}%
  \fancyhead[R]{\footnotesize\color{coulPart}\textbf{#2}}%
  \fancyfoot[C]{\footnotesize\lycee@pied}%
  \renewcommand{\headrulewidth}{0.4pt}%
  \renewcommand{\headrule}{\hbox to\headwidth{%
    \color{coulPart!40}\leaders\hrule height \headrulewidth\hfill}}}

\RequirePackage{ccicons}
\newcommand{\lycee@licence}{{\scriptsize\color{gray}%
  \copyright\ Luc Destombe --- \ccbyncsa\ CC BY-NC-SA 4.0}}
\newif\iflycee@licence \lycee@licencetrue
\newcommand{\sanslicence}{\lycee@licencefalse}
\AtBeginDocument{\iflycee@licence\AddToHookNext{shipout/foreground}{%
  \put(0,0){\rlap{\kern\dimexpr1in+\hoffset+\oddsidemargin+\textwidth\relax
    \llap{\raisebox{-\dimexpr1in+\voffset+\topmargin+\headheight+\headsep
      +\textheight+\footskip\relax}{\lycee@licence}}}}}\fi}

\newcommand{\formulesserrees}{%
  \setlength{\abovedisplayskip}{3pt}\setlength{\belowdisplayskip}{3pt}%
  \setlength{\abovedisplayshortskip}{2pt}\setlength{\belowdisplayshortskip}{3pt}}

\setlength{\parindent}{0pt}

\RequirePackage[implicit=false,hidelinks,pdfencoding=auto,psdextra,bookmarksopen,
  bookmarksopenlevel=1,pdfcreator={},pdfproducer={}]{hyperref}
\RequirePackage{bookmark}
\hypersetup{pdfauthor={Luc Destombe},pdfkeywords={CC BY-NC-SA 4.0}}
\newcounter{lycee@signet}
\newcommand{\signet}[3][]{\stepcounter{lycee@signet}%
  \edef\lycee@dest{\if\relax\detokenize{#1}\relax
    signet.\arabic{lycee@signet}\else#1\fi}%
  \hypertarget{\lycee@dest}{}\bookmark[level=#2,dest=\lycee@dest]{#3}}
\pdfstringdefDisableCommands{%
  \def\R{R}\def\vect#1{#1}\def\Cf{Cf}\def\Cg{Cg}\def\fois#1{×#1}%
  \def\barre#1{#1}\def\intervalle#1#2{[#1;#2]}\def\intff#1#2{[#1;#2]}%
  \def\textsuperscript#1{#1}\def\dfrac#1#2{#1/#2}\def\frac#1#2{#1/#2}%
  \def\vec#1{#1⃗}}%
\RequirePackage[official]{eurosym}

\tcbset{exercice corrige/.style={
  enhanced,breakable,before upper=\raggedright,
  colback=white,colframe=coulExo,
  boxrule=0.8pt,arc=2pt,
  left=6pt,right=6pt,top=8pt,bottom=5pt,
  before skip=9pt,after skip=4pt,
  coltitle=white,fonttitle=\bfseries\small,
  attach boxed title to top left={xshift=8pt,yshift=-\tcboxedtitleheight/2},
  boxed title style={colback=coulExo,arc=2pt,boxrule=0pt}}}
\newcounter{exo}
\newtcolorbox{exercice}[1][]{exercice corrige,
  phantom={\signet[exercice-\arabic{exo}]{2}{Exercice \arabic{exo}}},
  title={Exercice \theexo\ifx\relax#1\relax\else{} --- #1\fi}}
\newenvironment{exo}[1][\relax]{\refstepcounter{exo}\begin{exercice}[#1]}{\end{exercice}}

\RequirePackage{cuted}
\geometry{top=1.6cm,bottom=1cm,left=1cm,right=1cm,headheight=14pt,footskip=0.6cm}

\newcommand{\entetecorrige}[3]{%
  \begin{strip}
  \begin{center}
    \begin{tcolorbox}[enhanced,width=0.92\linewidth,colback=white,colframe=coulPart,
        boxrule=1.6pt,arc=3pt,halign=center,top=5pt,bottom=5pt]
      {\LARGE\bfseries\color{coulPart} #1}\\[3pt]
      {\normalsize #2}
    \end{tcolorbox}
  \end{center}
  \if\relax\detokenize{#3}\relax\else

  \vspace{2pt}
  \noindent
  \begin{tcolorbox}[enhanced,colback=white,colframe=black!35,boxrule=0.5pt,arc=2pt,
      left=7pt,right=7pt,top=4pt,bottom=4pt]
  {\small #3}
  \end{tcolorbox}
  \vspace{6pt}
  \fi
  \end{strip}}

\newcounter{partie}
\newcommand{\partiebandeau}[1]{%
  \besoinplace{8\baselineskip}\vspace{6pt}%
  \begin{tcolorbox}[enhanced,colback=coulPart!8,colframe=coulPart,boxrule=0pt,
      leftrule=3pt,arc=0pt,left=5pt,right=4pt,top=3pt,bottom=3pt,
      before skip=4pt,after skip=2pt]
    \bfseries\color{coulPart}#1\signet{1}{#1}
  \end{tcolorbox}}
\newcommand{\partie}[1]{\stepcounter{partie}\partiebandeau{\Roman{partie}) #1}}
\newcommand{\partiesansnum}[1]{\partiebandeau{#1}}

\newcommand{\res}[1]{%
  \tcbox[on line,colback=coulPart!7,colframe=coulPart,boxrule=0.5pt,arc=1pt,
    boxsep=0pt,left=3pt,right=3pt,top=1pt,bottom=2pt]{$#1$}}
\newcommand{\calc}[1]{\par\smallskip\noindent$\displaystyle #1$\par}
\newenvironment{calculs}{\par\smallskip\noindent$\displaystyle\begin{aligned}[t]}%
  {\end{aligned}$\par\smallskip}
\newcommand{\rem}[1]{\par\smallskip{\small\itshape\color{coulPart}#1}\par}
\newcommand{\deux}[2]{\par\noindent\makebox[0.5\linewidth][l]{#1}#2}

\setlist[questions,1]{label=\arabic*\textdegree),leftmargin=*,itemsep=2pt,topsep=2pt}
\setlist[questions,2]{label=\alph*),leftmargin=*,itemsep=2pt,topsep=1pt}
\newlist{reponses}{enumerate}{1}
\setlist[reponses]{label=\alph*),leftmargin=*,itemsep=2pt,topsep=2pt}

\setlength{\parskip}{1pt}
\setlength{\columnsep}{0.6cm}
\raggedbottom
\formulesserrees
\AtBeginDocument{\formulesserrees}

\makeatother
\usepackage{tkz-tab}
\entete{Ch 2 --- Fonctions affines et pourcentages --- Corrigé des exercices}{1\textsuperscript{re} mathématiques spécifiques}
\providecommand{\justif}[1]{\quad\text{\small\itshape\color{coulPart}#1}}

\begin{document}

\entetecorrige{FONCTIONS AFFINES ET POURCENTAGES}{Corrigé --- Ch 2 : Fiche d'exercices --- 1\textsuperscript{re} mathématiques spécifiques}
{Les remarques en italique signalent les erreurs les plus fréquentes et les méthodes à
retenir. Les résultats sont encadrés. Tous les calculs se font sans calculatrice.}

\partie{Appliquer un pourcentage}

\begin{exo}[Calcul mental]
\deux{a) \res{42}}{b) \res{15} \quad (le quart)}
\deux{c) \res{24} \quad (la moitié)}{d) \res{9}}
\deux{e) \res{4} \quad (moitié de $10\,\%$)}{f) \res{15}}
\deux{g) \res{7} \quad ($2\times 3{,}5$)}{h) \res{75} \quad ($3\times 25$)}
\deux{i) \res{6} \quad ($4\times 1{,}5$)}{}
\rem{On passe par $10\,\%$ (diviser par $10$) ou $1\,\%$ (diviser par $100$).}
\end{exo}

\begin{exo}[Trois écritures d'un même nombre]
\begin{center}
\renewcommand{\arraystretch}{1.8}
\small
\begin{tabular}{|l|c|c|c|c|c|}
\hline
Fraction & $\dfrac{3}{5}$ & \res{\dfrac{9}{100}} & \res{\dfrac{9}{20}} & $\dfrac{1}{4}$ & \res{\dfrac{3}{2}}\\
\hline
Décimal & \res{0{,}6} & $0{,}09$ & \res{0{,}45} & \res{0{,}25} & $1{,}5$\\
\hline
Pourc. & \res{60\,\%} & \res{9\,\%} & $45\,\%$ & \res{25\,\%} & \res{150\,\%}\\
\hline
\end{tabular}
\end{center}
\rem{$0{,}09=9\,\%$, et non $90\,\%$ ; un pourcentage peut dépasser $100\,\%$.}
\end{exo}

\begin{exo}[Retrouver le total]
\begin{reponses}
  \item $5\,\%$ est la vingtième partie : $20\times 30=\res{600}$ animaux.
  \item $30\,\%$ font $9$, donc $10\,\%$ font $3$ : \res{30} élèves.
  \item $10\,\%$ font $8$~€, donc $100\,\%$ font \res{80\text{ €}}.
  \item $75\,\%$ font $60$, donc $25\,\%$ font $20$ : \res{80} adhérents.
\end{reponses}
\end{exo}

\begin{exo}[Augmenter ou diminuer]
\begin{reponses}
  \item Hausse : $15\,\%$ de $60$ font $9$ ; nouveau prix \res{69\text{ €}}.
  \item Baisse : $20\,\%$ de $900$ font $180$ ; prix soldé \res{720\text{ €}}.
  \item Hausse : $3\,\%$ de $2\,000$ font $60$ ; salaire \res{2\,060\text{ €}}.
  \item Baisse : $5\,\%$ de $6\,000$ font $300$ ; \res{5\,700} habitants.
\end{reponses}
\end{exo}

\begin{exo}[Prix hors taxe, prix TTC]
\begin{reponses}
  \item TVA : $20\,\%$ de $350$ font \res{70\text{ €}} ; prix TTC
        $350+70=\res{420\text{ €}}$.
  \item $10\,\%$ de $420$ font $42$ : prix soldé \res{378\text{ €}}.
\end{reponses}
\rem{La baisse porte sur le prix TTC, $420$~€, et non sur le prix HT.}
\end{exo}

\partie{Calculer un pourcentage}

\begin{exo}[Proportions]
\begin{reponses}
  \item $\dfrac{12}{40}=\dfrac{30}{100}$, soit \res{30\,\%}
  \item $\dfrac{9}{20}=\dfrac{45}{100}$, soit \res{45\,\%}
  \item $\dfrac{33}{150}=\dfrac{22}{100}$, soit \res{22\,\%}
  \item $\dfrac{80}{500}=\dfrac{16}{100}$, soit \res{16\,\%}
\end{reponses}
\end{exo}

\begin{exo}[Deux enquêtes]
\begin{reponses}
  \item $\dfrac{140}{250}=\dfrac{56}{100}$, soit $56\,\%<60\,\%$ : \res{\text{non}}.
  \item $\dfrac{340}{400}=\dfrac{85}{100}$, soit $85\,\%$ : \res{\text{c'est exact}}.
\end{reponses}
\rem{Autre rédaction pour a) : $60\,\%$ de $250$ font $150$, et $140<150$.}
\end{exo}

\begin{exo}[Taux d'évolution]
\begin{reponses}
  \item $t=\dfrac{180-150}{150}=\dfrac{30}{150}$, soit $t=0{,}2$ : \res{+20\,\%}.
  \item $t=\dfrac{340-400}{400}=\dfrac{-60}{400}$, soit $t=-0{,}15$ : \res{-15\,\%}.
  \item $t=\dfrac{720-800}{800}=\dfrac{-80}{800}$, soit $t=-0{,}1$ : \res{-10\,\%}.
  \item $t=\dfrac{75-60}{60}=\dfrac{15}{60}$, soit $t=0{,}25$ : \res{+25\,\%}.
\end{reponses}
\rem{On divise toujours par la valeur de départ.}
\end{exo}

\begin{exo}[Plus ou moins de\ldots{} ?]
\begin{reponses}
  \item Perte : $80$ ; $10\,\%$ de $700$ font $70$ ; $80>70$ : \res{\text{oui}}.
  \item Gain : $120$ ; $10\,\%$ de $1\,500$ font $150$ ; $120<150$ : \res{\text{non}}.
  \item Perte : $120$ ; $5\,\%$ de $3\,000$ font $150$ ; $120<150$ : \res{\text{non}}.
\end{reponses}
\rem{Autre rédaction, celle des corrigés de l'épreuve, pour a) : une baisse de
$10\,\%$ donnerait $700\times 0{,}9=630$ ; or $620<630$.}
\end{exo}

\begin{exo}[Au tableur]
\begin{reponses}
  \item $t=\dfrac{5\,500-5\,000}{5\,000}=\dfrac{500}{5\,000}$, soit $t=0{,}1$ :
        \res{+10\,\%}.
  \item \res{\texttt{=(C2-B2)/B2}}
  \item \res{\texttt{=(D2-C2)/C2}}
\end{reponses}
\rem{Sans parenthèses, \texttt{=C2-B2/B2} calcule $C2-1$ : la division passe avant la
soustraction. On divise par la cellule de l'année \emph{précédente}.}
\end{exo}

\partie{Coefficient multiplicateur}

\begin{exo}[Du taux au coefficient]
\deux{a) \res{1{,}06} \quad b) \res{0{,}94}}{c) \res{1{,}3} \quad d) \res{0{,}7}}
\deux{e) \res{1{,}65} \quad f) \res{0{,}88}}{g) \res{1{,}04} \quad h) \res{0{,}55}}
\deux{i) \res{3}}{}
\rem{Hausse : $1+\frac{t}{100}$ ; baisse : $1-\frac{t}{100}$. Une hausse de $200\,\%$
triple la valeur.}
\end{exo}

\begin{exo}[Du coefficient au taux]
\deux{a) hausse de \res{9\,\%}}{b) baisse de \res{9\,\%}}
\deux{c) $\frac{4}{5}=0{,}8$ : baisse de \res{20\,\%}}{d) baisse de \res{30\,\%}}
\deux{e) hausse de \res{200\,\%}}{f) baisse de \res{2\,\%}}
\end{exo}

\begin{exo}[Appliquer une évolution]
\begin{reponses}
  \item $300\times 1{,}15=\res{345\text{ €}}$
  \item $600\times 0{,}65=\res{390\text{ €}}$
  \item $50\times 2{,}2=\res{110\text{ €}}$
  \item $2\,000\times 0{,}985=\res{1\,970}$ habitants
\end{reponses}
\rem{c) $+120\,\%$ : $C_{M}=1+1{,}2=2{,}2$.}
\end{exo}

\begin{exo}[Le bon calcul]
\begin{questions}
  \item \res{\text{b) } 500\times 0{,}88} : $C_{M}=1-0{,}12$.
  \item \res{\text{c) } 150\times 1{,}3} : $C_{M}=1+0{,}3$.
\end{questions}
\rem{$500\times 0{,}12$ est le montant de la baisse, pas le nouveau prix.}
\end{exo}

\begin{exo}[Retrouver un taux avec le coefficient]
\begin{reponses}
  \item $C_{M}=\dfrac{100}{80}=1{,}25$ : \res{+25\,\%}
  \item $C_{M}=\dfrac{400}{500}=0{,}8$ : \res{-20\,\%}
  \item $C_{M}=\dfrac{60}{20}=3$ : \res{+200\,\%}
\end{reponses}
\end{exo}

\partie{Évolutions successives, évolution réciproque}

\begin{exo}[Taux global]
\begin{reponses}
  \item $1{,}5\times 1{,}2=1{,}8$ : \res{+80\,\%}
  \item $0{,}7\times 1{,}3=0{,}91$ : \res{-9\,\%}
  \item $2\times 0{,}5=1$ : \res{0\,\%} (retour au départ)
  \item $0{,}8\times 0{,}8=0{,}64$ : \res{-36\,\%}
  \item $1{,}6\times 0{,}75=1{,}2$ : \res{+20\,\%}
\end{reponses}
\rem{Les taux ne s'ajoutent pas : on multiplie les coefficients.}
\end{exo}

\begin{exo}[Le prix de l'électricité]
\begin{calculs}
C_{M} &= 1{,}1\times 1{,}1\\
C_{M} &= 1{,}21
\end{calculs}
La hausse est de \res{21\,\%} et non de $20\,\%$ : le journaliste \res{\text{a tort}}.
\end{exo}

\begin{exo}[Vrai ou faux ?]
\begin{reponses}
  \item \res{\text{Faux}} : $1{,}5\times 0{,}5=0{,}75$, baisse de $25\,\%$.
  \item \res{\text{Vrai}} : $0{,}8\times 0{,}7=0{,}7\times 0{,}8=0{,}56$.
  \item \res{\text{Vrai}} : $0{,}9\times 0{,}9=0{,}81$.
\end{reponses}
\end{exo}

\begin{exo}[La valeur d'un scooter]
\begin{reponses}
  \item $4\,000\times 0{,}75=\res{3\,000\text{ €}}$ ; \;
        $3\,000\times 0{,}8=\res{2\,400\text{ €}}$.
  \item $C_{M}=0{,}75\times 0{,}8=0{,}6$ : \res{-40\,\%}.
\end{reponses}
\rem{Et non $-45\,\%$ : la seconde baisse porte sur $3\,000$~€.}
\end{exo}

\begin{exo}[Une même évolution répétée]
\begin{reponses}
  \item Multipliée par $1{,}03^{3}\approx 1{,}09$ : hausse d'environ \res{9\,\%}.
  \item $C_{M}=1{,}01^{12}\approx 1{,}13$ : hausse d'environ \res{13\,\%} (et non
        $12\,\%$).
\end{reponses}
\end{exo}

\begin{exo}[Se compensent-elles ?]
\deux{a) $0{,}8\times 1{,}25=1$ : \res{\text{oui}}}{b) $1{,}5\times 0{,}5=0{,}75$ : \res{\text{non}}}
\deux{c) $0{,}5\times 2=1$ : \res{\text{oui}}}{d) $4\times 0{,}25=1$ : \res{\text{oui}}}
\deux{e) $0{,}8\times 1{,}2=0{,}96$ : \res{\text{non}}}{f) $1{,}1\times 0{,}9=0{,}99$ : \res{\text{non}}}
\rem{Une hausse de $t\,\%$ n'est jamais compensée par une baisse de $t\,\%$.}
\end{exo}

\begin{exo}[Quelle évolution compense ?]
\begin{questions}
  \item \res{\text{b}} : $4\times 0{,}25=1$.
  \item \res{\text{c}} : $1{,}25\times 0{,}8=1$ (avec $25\,\%$ : $0{,}9375$).
  \item \res{\text{a}} : $0{,}5\times 2=1$ (avec $50\,\%$ : $0{,}75$).
\end{questions}
\end{exo}

\begin{exo}[Un verger]
\begin{reponses}
  \item $2\,000\times 0{,}8=\res{1\,600}$ arbres.
  \item $1\,600\times 1{,}2=1\,920$ : \res{\text{non}}.
  \item $1\,600\times 1{,}25=\res{2\,000}$ ; ou $0{,}8\times 1{,}25=1$.
\end{reponses}
\end{exo}

\partie{Fonctions affines : image, antécédent}

\begin{exo}[Affine ou non ?]
\begin{reponses}
  \item \res{\text{affine}}, $a=5$ et $b=1$
  \item \res{\text{non}} (carré de $x$)
  \item \res{\text{affine}}, $a=-3$ et $b=2$
  \item \res{\text{non}} ($x$ au dénominateur)
  \item \res{\text{affine}} (linéaire), $a=1{,}5$ et $b=0$
  \item \res{\text{affine}} (constante), $a=0$ et $b=-4$
\end{reponses}
\rem{$2-3x=-3x+2$ : le coefficient de $x$ est $a$, même s'il est écrit en second.}
\end{exo}

\begin{exo}[Images]
\begin{reponses}
  \item $f(0)=\res{-4}$ ; $f(2)=\res{2}$ ; $f(-1)=\res{-7}$ ; $f(0{,}5)=\res{-2{,}5}$
  \item $g(0)=\res{6}$ ; $g(2)=\res{-2}$ ; $g(-3)=\res{18}$ ; $g(1{,}5)=\res{0}$
\end{reponses}
\rem{$g(-3)=-4\times(-3)+6=12+6$ : parenthèses autour de $-3$.}
\end{exo}

\begin{exo}[Antécédents]
a) Antécédent de $15$ :
\begin{calculs}
3x+9 &= 15\\
3x &= 6\\
x &= \res{2}
\end{calculs}
De même : $3x=-9$, antécédent de $0$ : \res{-3} ; \; $3x=-15$, antécédent de $-6$ :
\res{-5}.

\smallskip
b) Antécédent de $14$ :
\begin{calculs}
-2x+8 &= 14\\
-2x &= 6\\
x &= \res{-3}
\end{calculs}
Antécédent de $0$ : $-2x=-8$, soit \res{4}.

\smallskip
c) $0{,}5x-2=3$, donc $0{,}5x=5$ et \res{x=10}.
\end{exo}

\begin{exo}[Deux tarifs]
a) \res{T(x)=1{,}5x+40}.
\begin{calculs}
1{,}5x+40 &= 70\\
1{,}5x &= 30\\
x &= \res{20} &&\justif{$\frac{30}{1{,}5}=\frac{300}{15}$}
\end{calculs}
Le dépanneur a parcouru $20$~km.

\smallskip
b) \res{S(x)=15x+200} ; \; $S(40)=600+200=\res{800\text{ €}}$.
\begin{calculs}
15x+200 &= 650\\
15x &= 450\\
x &= \res{30} \text{ invités}
\end{calculs}
\end{exo}

\partie{Trouver une fonction affine}

\begin{exo}[Mettre en équation]
\begin{reponses}
  \item \res{L(x)=6x+8}
  \item $5$~L pour $100$~km, soit $0{,}05$~L par km : \res{R(x)=60-0{,}05x}
  \item $4\,\%$ des ventes : $0{,}04x$ ; \; \res{S(x)=0{,}04x+1\,400}
\end{reponses}
\rem{$b$ est la valeur fixe ; $a$ ce que coûte ou rapporte une unité de $x$.}
\end{exo}

\begin{exo}[Coefficient directeur]
\begin{reponses}
  \item $a=\dfrac{11-2}{4-1}=\dfrac{9}{3}$, soit \res{a=3}
  \item $a=\dfrac{2-6}{2-0}=\dfrac{-4}{2}$, soit \res{a=-2}
  \item $a=\dfrac{7-3}{-4-(-2)}=\dfrac{4}{-2}$, soit \res{a=-2}
  \item $a=\dfrac{6-(-2)}{1-(-3)}=\dfrac{8}{4}$, soit \res{a=2}
\end{reponses}
\rem{Parenthèses autour des coordonnées négatives : $-4-(-2)=-4+2=-2$.}
\end{exo}

\begin{exo}[Déterminer une fonction affine]
a) $a=\dfrac{11-5}{3-1}=3$. Avec $A$ :
\begin{calculs}
3\times 1+b &= 5\\
b &= 2
\end{calculs}
\res{f(x)=3x+2}.

\smallskip
b) $a=\dfrac{2-(-4)}{2-0}=3$ et $b=-4$ (ordonnée de $A$) : \res{f(x)=3x-4}.

\smallskip
c) $a=\dfrac{3-3}{2-(-2)}=0$ : \res{f(x)=3} (fonction constante).

\smallskip
d) $a=\dfrac{4-(-5)}{2-(-4)}=\dfrac{9}{6}$, soit $a=1{,}5$. Avec $B$ :
\begin{calculs}
1{,}5\times 2+b &= 4\\
3+b &= 4\\
b &= 1
\end{calculs}
\res{f(x)=1{,}5x+1}.
\rem{Contrôle avec l'autre point : $f(-4)=-6+1=-5$.}
\end{exo}

\begin{exo}[Le dioxyde de carbone dans l'air]
\begin{reponses}
  \item $f(0)=\res{370}$ et $f(10)=\res{390}$ ; $a=\dfrac{390-370}{10-0}=2$ et
        $b=370$ : \res{f(x)=2x+370}.
  \item La concentration augmente de \res{2\text{ ppm par an}}.
  \item 2030 : $x=30$ ; $f(30)=60+370=\res{430\text{ ppm}}$.
  \item $2x+370=450$, donc $2x=80$ et $x=40$ : en \res{2040}.
\end{reponses}
\end{exo}

\partie{Tracer une fonction affine}

\begin{exo}[Sens de variation]
\begin{reponses}
  \item $a=-1<0$ : \res{\text{décroissante}}
  \item $a=3>0$ : \res{\text{croissante}}
  \item $a=0$ : \res{\text{constante}}
  \item $a=-5<0$ : \res{\text{décroissante}}
  \item $a=-0{,}2<0$ : \res{\text{décroissante}}
  \item $a=0{,}3>0$ : \res{\text{croissante}}
\end{reponses}
\rem{Seul le signe de $a$ compte, pas celui de $b$.}
\end{exo}

\begin{exo}[Tracer des droites]
\begin{minipage}[t]{0.46\linewidth}\raggedright\vspace{0pt}
$f$ : $(0\,;1)$ et $(2\,;3)$.

$g$ : $(0\,;2)$ et $(1\,;-1)$.

$h$ : $(0\,;-2)$ et $(4\,;0)$.

\smallskip
{\small\itshape\color{coulPart}Pour $h$, on choisit une abscisse paire : $h(4)$ est
entier.}
\end{minipage}\hfill
\begin{minipage}[t]{0.52\linewidth}\centering\vspace{0pt}
\begin{tikzpicture}[x=0.36cm,y=0.36cm]
  \quadrillage{-3}{-3}{5}{4}
  \draw[->,>=Stealth,black!70] (-3,0) -- (5.5,0) node[below,font=\tiny] {$x$};
  \draw[->,>=Stealth,black!70] (0,-3) -- (0,4.5) node[left,font=\tiny] {$y$};
  \foreach \i in {-2,-1,1,2,3,4} {\node[below,font=\tiny,inner sep=1.5pt] at (\i,0) {$\i$};}
  \foreach \i in {-2,-1,1,2,3} {\node[left,font=\tiny,inner sep=1.5pt] at (0,\i) {$\i$};}
  \begin{scope}
    \clip (-3,-3) rectangle (5,4);
    \draw[coulCourbe,thick] (-3,-2) -- (3,4);
    \draw[coulCourbeB,thick] (-0.667,4) -- (1.667,-3);
    \draw[coulCourbeC,thick] (-2,-3) -- (5,0.5);
  \end{scope}
  \foreach \p in {(0,1),(2,3),(0,2),(1,-1),(0,-2),(4,0)} {\path \p node[croix]{};}
  \node[coulCourbe,font=\small,left] at (2.6,3.7) {$f$};
  \node[coulCourbeB,font=\small,right] at (1.5,-2.6) {$g$};
  \node[coulCourbeC,font=\small,above] at (4.4,0.3) {$h$};
\end{tikzpicture}
\end{minipage}
\end{exo}

\begin{exo}[Location de paddles]
a) \res{L(x)=4x+4} \qquad b) $L(0)=\res{4}$ ; \; $L(8)=\res{36}$.

\smallskip
\begin{minipage}[t]{0.44\linewidth}\raggedright\vspace{0pt}
c) Droite passant par $(0\,;4)$ et $(8\,;36)$ (figure réduite).

\smallskip
d) Pour $24$~€, on lit \res{5\text{ heures}}. Calcul : $4x+4=24$, donc $4x=20$ et
$x=5$.
\end{minipage}\hfill
\begin{minipage}[t]{0.54\linewidth}\centering\vspace{0pt}
\begin{tikzpicture}[x=0.42cm,y=0.07cm]
  \draw[grille,xstep=1,ystep=5] (0,0) grid (8,40);
  \draw[->,>=Stealth,black!70] (0,0) -- (8.6,0) node[below,font=\tiny] {h};
  \draw[->,>=Stealth,black!70] (0,0) -- (0,43) node[left,font=\tiny] {€};
  \foreach \i in {1,...,8} {\node[below,font=\tiny,inner sep=1.5pt] at (\i,0) {$\i$};}
  \foreach \v in {5,10,...,40} {\node[left,font=\tiny,inner sep=1.5pt] at (0,\v) {$\v$};}
  \draw[coulCourbe,thick] (0,4) -- (8,36);
  \draw[dashed,coulRep] (0,24) -- (5,24) -- (5,0);
  \foreach \p in {(0,4),(8,36)} {\path \p node[croix]{};}
\end{tikzpicture}
\end{minipage}
\rem{Repère non orthonormé : les deux grandeurs n'ont pas la même unité.}
\end{exo}

\begin{exo}[Proportionnalité et fonctions linéaires]
\begin{reponses}
  \item $a=\dfrac{10}{100}=0{,}1$ : \res{f(x)=0{,}1x} ; \; $f(330)=\res{33\text{ g}}$.
  \item $40$ pouces : $2{,}5\times 40=\res{100\text{ cm}}$ ; \; $2{,}5x=75$ donne
        \res{x=30} pouces.
  \item \res{f} et \res{h} ($h(x)=0{,}25x$) ; $g$ est affine non linéaire, $k$ est
        constante.
  \item $a=\dfrac{15}{5}=3$ : \res{\ell(x)=3x}.
\end{reponses}
\rem{Linéaire : $b=0$, la droite passe par l'origine.}
\end{exo}

\partie{Lire l'équation d'une droite}

\begin{exo}[Lire une équation réduite]
\begin{reponses}
  \item[$(d_{1})$ :] $b=-1$ ; de $(0\,;-1)$ à $(1\,;2)$ : $a=\dfrac{3}{1}=3$ ;
        \res{y=3x-1}, croissante.
  \item[$(d_{2})$ :] $b=2$ ; de $(0\,;2)$ à $(2\,;0)$ : $a=\dfrac{-2}{2}=-1$ ;
        \res{y=-x+2}, décroissante.
  \item[$(d_{3})$ :] $b=1$ ; de $(0\,;1)$ à $(3\,;2)$ : $a=\dfrac{1}{3}$ ;
        \res{y=\tfrac{1}{3}x+1}, croissante.
  \item[$(d_{4})$ :] horizontale : \res{y=-2}, constante.
\end{reponses}
\rem{$a=\dfrac{\text{déplacement vertical}}{\text{déplacement horizontal}}$ : pour
$(d_{3})$, on avance de $3$ pour monter de $1$.}
\end{exo}

\begin{exo}[Points d'une droite]
\begin{reponses}
  \item $A$ : $-2\times 1+5=3$ : \res{\text{oui}} ; \; $B$ : $-2\times 3+5=-1$ et
        $-1\ne 1$ : \res{\text{non}} ; \; $C$ : $-2\times(-2)+5=9$ : \res{\text{oui}}.
  \item $y=-2\times 4+5=\res{-3}$.
  \item $-2x+5=-7$, donc $-2x=-12$ et \res{x=6}.
\end{reponses}
\end{exo}

\begin{exo}[Équation réduite sans graphique]
\begin{questions}
  \item \res{\text{b}} : $b=2$ et $a=\dfrac{1}{3}$.
  \item \res{\text{c}} : $b=-3$ et $a=\dfrac{1-(-3)}{2-0}=2$.
  \item \res{\text{a}} : $a=4>0$.
\end{questions}
\end{exo}

\partie{Signe d'une fonction affine}

\begin{exo}[Tableaux de signes]
a) $3x-9=0$, \res{x=3} ; $a>0$\hfill
\begin{tikzpicture}[baseline=(current bounding box.center)]
\tkzTabInit[lgt=0.9,espcl=1.1,deltacl=0.3]{$x$/0.55,$f(x)$/0.55}{$-\infty$,$3$,$+\infty$}
\tkzTabLine{,-,z,+,}
\end{tikzpicture}\par\smallskip
b) $-2x+8=0$, \res{x=4} ; $a<0$\hfill
\begin{tikzpicture}[baseline=(current bounding box.center)]
\tkzTabInit[lgt=0.9,espcl=1.1,deltacl=0.3]{$x$/0.55,$g(x)$/0.55}{$-\infty$,$4$,$+\infty$}
\tkzTabLine{,+,z,-,}
\end{tikzpicture}\par\smallskip
c) $4x+12=0$, \res{x=-3} ; $a>0$\hfill
\begin{tikzpicture}[baseline=(current bounding box.center)]
\tkzTabInit[lgt=0.9,espcl=1.1,deltacl=0.3]{$x$/0.55,$h(x)$/0.55}{$-\infty$,$-3$,$+\infty$}
\tkzTabLine{,-,z,+,}
\end{tikzpicture}\par\smallskip
d) $-0{,}5x+3=0$, \res{x=6} ; $a<0$\hfill
\begin{tikzpicture}[baseline=(current bounding box.center)]
\tkzTabInit[lgt=0.9,espcl=1.1,deltacl=0.3]{$x$/0.55,$k(x)$/0.55}{$-\infty$,$6$,$+\infty$}
\tkzTabLine{,+,z,-,}
\end{tikzpicture}\par\smallskip
\rem{À droite du zéro, le signe de $a$ ; à gauche, le signe contraire.}
\end{exo}

\begin{exo}[Positif ou négatif ?]
\begin{questions}
  \item \res{\text{a}} : $-5x+2=0$ pour $x=0{,}4$ et $a<0$ : positif avant $0{,}4$.
  \item \res{\text{c}} : $3x+12=0$ pour $x=-4$ et $a>0$ : négatif avant $-4$.
\end{questions}
\end{exo}

\begin{exo}[Un bénéfice]
\begin{reponses}
  \item $B(0)=\res{-150}$ : sans vente, le food-truck perd $150$~€ (ses frais fixes).
  \item $5x-150=0$, donc $5x=150$ et \res{x=30}.
  \item $a=5>0$ :\par
\begin{tikzpicture}
\tkzTabInit[lgt=1.2,espcl=1.3]{$x$/0.6,$B(x)$/0.6}{$0$,$30$,$80$}
\tkzTabLine{,-,z,+,}
\end{tikzpicture}
  \item Le bénéfice est positif à partir de \res{31} repas (nul pour $30$).
\end{reponses}
\end{exo}

\partie{Comparer deux fonctions affines}

\begin{exo}[Équations]
a)
\begin{calculs}
6x+5 &= 4x+11\\
2x &= 6\\
x &= \res{3}
\end{calculs}
b)
\begin{calculs}
3x-2 &= x+8\\
2x &= 10\\
x &= \res{5}
\end{calculs}
c)
\begin{calculs}
2x+3 &= 5x-9\\
-3x &= -12\\
x &= \res{4}
\end{calculs}
d)
\begin{calculs}
4x+1 &= -2x+4\\
6x &= 3\\
x &= \res{0{,}5}
\end{calculs}
\rem{En d), on ajoute $2x$ des deux côtés : $4x+2x=6x$.}
\end{exo}

\begin{exo}[Deux offres d'emploi]
a) \res{f(x)=0{,}1x+1\,200} et \res{g(x)=0{,}2x+800}.

b)
\begin{calculs}
0{,}1x+1\,200 &= 0{,}2x+800\\
400 &= 0{,}1x\\
x &= \res{4\,000\text{ €}}
\end{calculs}
c) $f(5\,000)=500+1\,200=1\,700$ et $g(5\,000)=1\,000+800=1\,800$ :
\res{\text{société B}}.
\rem{En dessous de $4\,000$~€ de ventes, la société A est plus intéressante.}
\end{exo}

\begin{exo}[Deux formules de piscine]
\begin{reponses}
  \item \res{f(x)=5x} et \res{g(x)=2x+24}.
  \item Intersection : \res{(8\,;40)}.
  \item $5x=2x+24$, donc $3x=24$ et $x=8$ ; $f(8)=40$.
  \item Moins de $8$ entrées : \res{\text{A}} ; plus de $8$ : \res{\text{B}} ; pour
        $8$ entrées, les deux coûtent $40$~€.
\end{reponses}
\end{exo}

\partiesansnum{Exercices type épreuve anticipée}

\begin{exo}[Automatismes (QCM)]
\begin{questions}
  \item \res{\text{a}} : $3\times 7=21$ ($10\,\%$ de $70$ font $7$).
  \item \res{\text{b}} : $100\times 1{,}2\times 0{,}8=100\times 0{,}96$.
  \item \res{\text{c}} : $\frac{1}{8}=0{,}125$ ($\frac{1}{4}=25\,\%=0{,}25$).
  \item \res{\text{b}} : $0{,}8\times 1{,}25=1$.
  \item \res{\text{a}} : $\dfrac{7-1}{4-2}=\dfrac{6}{2}$.
  \item \res{\text{a}} : $a=3>0$.
  \item \res{\text{b}} : $2x+3=13$ donne $2x=10$.
  \item \res{\text{a}} : $-6+4=-2$.
  \item \res{\text{a}} : $1{,}6\times 25=40$.
\end{questions}
\end{exo}

\begin{exo}[Les inscrits d'un club de judo]
\begin{questions}
  \item $t=\dfrac{171-180}{180}=\dfrac{-9}{180}$, soit $t=-0{,}05$ : \res{-5\,\%}.
  \item \res{\texttt{=(C3-B3)/B3}}
  \item $a=-8<0$ : $f$ est \res{\text{décroissante}} ; chaque année, le club compte
        \res{8} inscrits de moins.
  \item 2028 : $x=8$ ; $f(8)=-64+180=\res{116}$ inscrits.
  \item On résout l'inéquation :
\begin{calculs}
-8x+180 &< 100\\
-8x &< -80\\
x &> 10 &&\justif{on divise par $-8$ : le sens change}
\end{calculs}
        Premier entier : $x=11$, soit \res{2031} ($f(10)=100$ n'est pas « moins de
        $100$ »).
\end{questions}
\end{exo}

\begin{exo}[Les soldes]
\begin{questions}
  \item Casque : $80\times 0{,}75=\res{60\text{ €}}$ ; enceinte :
        $120\times 0{,}9=\res{108\text{ €}}$.
  \item Total : $60+108=\res{168\text{ €}}$ au lieu de $200$~€ ;
        $t=\dfrac{168-200}{200}=\dfrac{-32}{200}$, soit une remise globale de
        \res{16\,\%}.
  \item $108\times 1{,}1=118{,}80$ : \res{\text{non}} ($0{,}9\times 1{,}1=0{,}99$
        et $0{,}99\ne 1$).
\end{questions}
\rem{La remise globale n'est ni $35\,\%$ ni la moyenne $17{,}5\,\%$ : l'enceinte, plus
chère, pèse davantage.}
\end{exo}

\end{document}
